% GENERATED FILE. Edit canonical lessons, then run npm run generate:books.
% canonical-source-hash: fnv1a64:215cd813972742cf
% canonical-lessons: SA-C151-utsahah, SA-C151-dhairyam, SA-C151-alasyam, SA-C151-matsarah, SA-C151-garvah

\chapter{Enthusiasm, Courage, Laziness, Jealousy, Pride}
\label{ch:sa-enthusiasm-courage-laziness-jealousy-pride}
\hlchaptermodality{151}

\begin{chapteropening}
\textbf{By the end of this chapter:} \emph{I can say enthusiasm, courage, laziness, jealousy and pride.}

Complete the last lesson of chapter 151: I can say enthusiasm, courage, laziness, jealousy and pride.
\end{chapteropening}

\section[\texorpdfstring{\sk{उत्साहः} (utsāhaḥ)}{उत्साहः (utsāhaḥ)}]{\sk{उत्साहः} (utsāhaḥ) — enthusiasm}

\label{lesson:SA-C151-utsahah}



\begin{quote}
Before the new one: say the Sanskrit for vomiting, then the Sanskrit for a stomach ache.
\end{quote}

\subsection*{You'll want to know: \sk{उत्साहः}}
\textbf{\sk{उत्साहः}} — \emph{utsāhaḥ} — \textquotedblleft{}enthusiasm\textquotedblright{}.

\begin{grammarlens}[title={What you've built}]
One more word for this chapter.
\end{grammarlens}

\subsection*{Your turn}
\begin{itemize}
\raggedright
  \item \emph{Say it:} \emph{utsāhaḥ}
  \item \emph{Say it:} \emph{utsāhaḥ}, once more
  \item \emph{Say it:} say \emph{utsāhaḥ}
  \item \emph{Recall:} say the Sanskrit for vomiting, then the Sanskrit for a stomach ache, then say \emph{utsāhaḥ} again
\end{itemize}

\begin{tcolorbox}[breakable,colback=teal!4,colframe=teal!35!black,title={Before you move on}]
What does \sk{उत्साहः} mean? (\textquotedblleft{}Enthusiasm\textquotedblright{}.) Say it once more.
\end{tcolorbox}
\section[\texorpdfstring{\sk{धैर्यम्} (dhairyam)}{धैर्यम् (dhairyam)}]{\sk{धैर्यम्} (dhairyam) — courage, patience}

\label{lesson:SA-C151-dhairyam}



\begin{quote}
Before the new one: say the Sanskrit for a stomach ache, then the Sanskrit for enthusiasm.
\end{quote}

\subsection*{You'll want to know: \sk{धैर्यम्}}
\textbf{\sk{धैर्यम्}} — \emph{dhairyam} — \textquotedblleft{}courage, patience\textquotedblright{}.

\begin{grammarlens}[title={What you've built}]
One more word for this chapter.
\end{grammarlens}

\subsection*{Your turn}
\begin{itemize}
\raggedright
  \item \emph{Say it:} \emph{dhairyam}
  \item \emph{Say it:} \emph{dhairyam}, once more
  \item \emph{Say it:} say \emph{dhairyam}
  \item \emph{Recall:} say the Sanskrit for a stomach ache, then the Sanskrit for enthusiasm, then say \emph{dhairyam} again
\end{itemize}

\begin{tcolorbox}[breakable,colback=teal!4,colframe=teal!35!black,title={Before you move on}]
What does \sk{धैर्यम्} mean? (\textquotedblleft{}Courage, patience\textquotedblright{}.) Say it once more.
\end{tcolorbox}
\section[\texorpdfstring{\sk{आलस्यम्} (ālasyam)}{आलस्यम् (ālasyam)}]{\sk{आलस्यम्} (ālasyam) — laziness}

\label{lesson:SA-C151-alasyam}



\begin{quote}
Before the new one: say the Sanskrit for enthusiasm, then the Sanskrit for courage.
\end{quote}

\subsection*{You'll want to know: \sk{आलस्यम्}}
\textbf{\sk{आलस्यम्}} — \emph{ālasyam} — \textquotedblleft{}laziness\textquotedblright{}.

\begin{grammarlens}[title={What you've built}]
One more word for this chapter.
\end{grammarlens}

\subsection*{Your turn}
\begin{itemize}
\raggedright
  \item \emph{Say it:} \emph{ālasyam}
  \item \emph{Say it:} \emph{ālasyam}, once more
  \item \emph{Say it:} say \emph{ālasyam}
  \item \emph{Recall:} say the Sanskrit for enthusiasm, then the Sanskrit for courage, then say \emph{ālasyam} again
\end{itemize}

\begin{tcolorbox}[breakable,colback=teal!4,colframe=teal!35!black,title={Before you move on}]
What does \sk{आलस्यम्} mean? (\textquotedblleft{}Laziness\textquotedblright{}.) Say it once more.
\end{tcolorbox}
\section[\texorpdfstring{\sk{मत्सरः} (matsaraḥ)}{मत्सरः (matsaraḥ)}]{\sk{मत्सरः} (matsaraḥ) — jealousy}

\label{lesson:SA-C151-matsarah}



\begin{quote}
Before the new one: say the Sanskrit for courage, then the Sanskrit for laziness.
\end{quote}

\subsection*{You'll want to know: \sk{मत्सरः}}
\textbf{\sk{मत्सरः}} — \emph{matsaraḥ} — \textquotedblleft{}jealousy\textquotedblright{}.

\begin{grammarlens}[title={What you've built}]
One more word for this chapter.
\end{grammarlens}

\subsection*{Your turn}
\begin{itemize}
\raggedright
  \item \emph{Say it:} \emph{matsaraḥ}
  \item \emph{Say it:} \emph{matsaraḥ}, once more
  \item \emph{Say it:} say \emph{matsaraḥ}
  \item \emph{Recall:} say the Sanskrit for courage, then the Sanskrit for laziness, then say \emph{matsaraḥ} again
\end{itemize}

\begin{tcolorbox}[breakable,colback=teal!4,colframe=teal!35!black,title={Before you move on}]
What does \sk{मत्सरः} mean? (\textquotedblleft{}Jealousy\textquotedblright{}.) Say it once more.
\end{tcolorbox}
\section[\texorpdfstring{\sk{गर्वः} (garvaḥ)}{गर्वः (garvaḥ)}]{\sk{गर्वः} (garvaḥ) — pride}

\label{lesson:SA-C151-garvah}



\begin{quote}
Before the new one: say the Sanskrit for laziness, then the Sanskrit for jealousy.
\end{quote}

\subsection*{You'll want to know: \sk{गर्वः}}
\textbf{\sk{गर्वः}} — \emph{garvaḥ} — \textquotedblleft{}pride\textquotedblright{}.

\begin{grammarlens}[title={What you've built}]
One more word for this chapter.
\end{grammarlens}

\subsection*{Your turn}
\begin{itemize}
\raggedright
  \item \emph{Say it:} \emph{garvaḥ}
  \item \emph{Say it:} \emph{garvaḥ}, once more
  \item \emph{Say it:} say \emph{garvaḥ}
  \item \emph{Recall:} say the Sanskrit for laziness, then the Sanskrit for jealousy, then say \emph{garvaḥ} again
\end{itemize}

\begin{tcolorbox}[breakable,colback=teal!4,colframe=teal!35!black,title={Before you move on}]
What does \sk{गर्वः} mean? (\textquotedblleft{}Pride\textquotedblright{}.) Say it once more.
\end{tcolorbox}
